\subsection{Motivation}

Suppose you need to develop the zero-knowledge circuit over some ``unfriendly''
construction such as AES-128, SHA-256, or pairing implemented on the non-native
finite field. One way to resolve this issue is to use the proving scheme that is
most ``native'' to the logic needed to implement such constructions (for
example, you might want to switch to STARK instead of Groth16 for implementing
hash functions). However, zero-knowledge world came up with one more way to
resolve this issue without changing the whole system.

Consider the following problem: you have the witness $\mathbbm{w}$ in which you
want to check whether elements, say, $\vec{z} := \{z_i\}_{i \in [n]} \subseteq
\mathbb{F}$ are contained in the lookup table $\vec{t} := \{t_j\}_{j \in [d]}
\subseteq \mathbb{F}$ of size $d$. We write this check as $\{z_i\}_{i \in [n]}
\subseteq \{t_j\}_{j \in [d]}$ or $\vec{z} \subseteq \vec{t}$ for short.

\begin{example}
    As a not very practical example, suppose $\vec{t} = \{0,1\}$. Then the check
    $\vec{z} \subseteq \vec{t}$ is read as ``check whether each of elements in
    $\vec{z}$ is binary''. Another trivial example might be the following: set
    $\vec{t} := \{1,6,7,10\}$. Then, for example, $\vec{z} =
    \{10,6,7,1,1,6,10,7,1\}$ should pass the lookup check since $\vec{z}$
    consists of elements from $\vec{t}$ only. In turn, $\vec{z}^*=\{1,6,10,5\}$
    is not a valid witness corresponding to $\vec{t}$ since $5 \not\in \vec{t}$.
\end{example}

Modern protocols have the following property: such check costs only
$\mathcal{O}(n+d)$ constraints. Let us consider why this is significant for many
applications in the wild.

\textbf{Range Checks.} The most obvious application is the range check. Assume
that $t_j=j$ and $d=2^w$. This way, the check $\{z_i\}_{i \in [n]} \subseteq
\{t_j\}_{j \in [d]}$ simply corresponds to checking whether ``is each $z_i$
indeed a $w$-bit integer''? Needless to say, this has significant implications
in various ZK applications: privacy-preserving payments, implementation of
non-native arithmetic, or even zkML. In the case of zkML, for instance,
\textit{Bionetta} uses lookup checks to speed up the computation of the ReLU
function, given by $\text{ReLU}(x) = \max\{0,x\}$, by the factor of up to
$\times 20$.

\textbf{Binary Arithmetic.} Suppose you are implementing the SHA256, in which
you need to frequently XOR binary strings. To simplify the problem, instead of
XORing two, say, 256-bit strings, you XOR 8-bit strings: $\mathbf{c} =
\mathbf{a} \oplus \mathbf{b}$ for $\mathbf{a},\mathbf{b},\mathbf{c} \in
\{0,1\}^8$. The naive approach is to bit decompose each of
$\mathbf{a},\mathbf{b},\mathbf{c}$ and verify that the check holds, but that
requires enormous number of constraints. The more clever lookup-based approach
is as follows: compose the table $\vec{t} := \{(t_j^{\langle 1
\rangle},t_j^{\langle 2 \rangle},t_j^{\langle 3 \rangle})\}_{j \in [d]}
\subseteq \mathbb{F}^3$ which consists of all distinct valid 8-bit XORs:
$t_j^{\langle 1 \rangle} \oplus t_j^{\langle 2 \rangle}=t_j^{\langle 3
\rangle}$. Note that in such case $d=2^{16}$ since there are $2^{16}$ distinct
equalities of such form (which simply enumerate over all possibilities of two
8-bit inputs to the XOR gate). We thus want to show that $\vec{z} :=
\{(a_i,b_i,c_i)\}_{i \in [n]} \subseteq \vec{t}$ where $a_i,b_i,c_i$ are the
part of the witness. Now what remains is to show how to reduce this check to the
``one-dimensional'' case: recall that lookup protocols can merely check for
inclusion for the ordered set of \textit{scalars} and not \textit{tuples}. It is
done fairly easy: sample some random $\alpha \leftarrowS \mathbb{F}$ and instead
of viewing $\vec{z}$ as a tuple of three elements, view it as a list
$\vec{z}_{\alpha}$ of $n$ scalars of form $a_i+\alpha b_i+\alpha^2 c_i$.
Similarly, each table elements of $\vec{t}_{\alpha}$ is viewed as $t_j^{\langle
1 \rangle} + \alpha t_j^{\langle 2 \rangle} + \alpha^2 t_j^{\langle 3 \rangle}$.
Then, the inclusion check $\vec{z} \subseteq \vec{t}$ due to Schwartz-Zippel
lemma is equivalent to $\vec{z}_{\alpha} \subseteq \vec{t}_{\alpha}$ with
overwhelming probability.

\subsection{Plookup Protocol}

One of the first lookup protocols that became practical in the zero-knowledge
world is the \textit{plookup protocol}. It is mostly used in Poly-IOPs but can
presumably be compiled to other types of protocols as well, as long as
\textit{multiset equality} check can be implemented optimally.

Now, if you recall the $\mathcal{P}$lon$\mathcal{K}$ construction, we can check
that the multiset $\{s_i'\}_{i \in [n]}$ is the permutation of the multiset
$\{s_i\}_{i \in [n]}$. This was done by comparing the products $\prod_{i \in
[n]}(\gamma+s_i')=\prod_{i \in [n]}(\gamma+s_i)$ at a random point $\gamma$. Due
to Schwartz-Zippel lemma, equality in such case implies $s_{i}=s'_{\sigma(i)}$
with overwhelming probability for aribtrary $i \in [n]$ and some permutation $\sigma
\in S_n$.

The primary issue with the lookup check is that instead of multisets equality,
we need to check whether one multiset is the subset of another, which is a much
trickier problem: we need to show that polynomials $Z(X) := \prod_{i \in
[n]}(X-z_i)$ and $T(X) := \prod_{j \in [d]}(X-t_j)$ have the same roots,
ignoring multiplicities. Let us introduce the following tool which will help to
reduce the problem to a simpler one.

\begin{definition}
    The \textbf{difference set} of $\vec{s} = \{s_i\}_{i \in [n]}$, denoted by
    $\delta\vec{s}$, is $\delta\vec{s} := \{s_{i+1}-s_i\}_{i \in [n-1]}$.
    Additionally, by $\partial\vec{s}$ denote the non-zero difference set: the
    set $\delta\vec{s}$, but without zero elements.
\end{definition}

\begin{example}
    Suppose $\vec{s} = \{1,2,5,10,13\}$. Then, we have 
    \begin{equation*}
        \delta\vec{s} = \{2-1,5-2,10-5,13-10\} = \{1,3,5,3\} = \partial\vec{s}
    \end{equation*}
    A slightly more motivating example is $\vec{s} := \{1,1,4,4,8,8,8\}$. Then:
    \begin{equation*}
        \delta\vec{s} = \{0,3,0,4,0,0\}, \quad \partial\vec{s} = \{3,4\}.
    \end{equation*}
\end{example}

Now let us get to the initial goal of proving $\vec{z} = \{z_i\}_{i \in [n]}
\subseteq \{t_j\}_{j \in [d]} =: \vec{t}$. Let us look at a sorted version
$\vec{s}$ of the values of $\vec{z}$ and compare the non-zero difference sets
$\partial\vec{s}$ and $\partial\vec{t}$. Notice the following: if $\vec{z}
\subseteq \vec{t}$, then $\partial\vec{s} = \partial\vec{t}$. So the prover
$\mathcal{P}$ might want to simply prove this set equality. Unfortunately, the
converse (that is, $\partial\vec{s} = \partial\vec{t} \Rightarrow \vec{z}
\subseteq \vec{t}$) is not generally true. Indeed, consider the following
example:
\begin{equation*}
    \vec{t} = \{1,4,8\}, \quad \vec{s} = \{1,1,4,8,8,8\}, \quad \vec{s}^{\,*} = \{1,5,5,5,8,8\}.
\end{equation*}

Note that $\partial\vec{t} = \partial\vec{s} = \partial\vec{s}^{\,*}=\{3,4\}$,
but obviously $\vec{s}^{\,*} \not\subseteq \vec{t}$. So we need to modify our
checks to introduce more robust lookup check condition. 

Now, consider the following version: let $\vec{s}$ be the \textit{sorted}
multiset of concatenation, which we denote by $(\vec{z}, \vec{t}) =
\{z_0,\dots,z_{n-1},t_0,\dots,t_{d-1}\}$. Now we impose two checks: (1) $\vec{s}
= (\vec{z}, \vec{t})$, (2) $\partial\vec{s} = \partial\vec{t}$. Turns out this
condition is necessary and sufficient.

\begin{proposition}
    $\vec{z} \subseteq \vec{t} \iff \vec{s} = (\vec{z}, \vec{t}) \wedge \partial\vec{s} =
    \partial\vec{t}$.
\end{proposition}

\textbf{Proof.} \textbf{$\Rightarrow$ direction.} The first condition holds due
to construction. Note that the second condition holds for the reason that
$\partial\vec{s}$ contains only differences of form $z_i-t_j$, $z_i-z_j$, or
$t_i-t_j$. If $\vec{z} \subseteq \vec{t}$ and given $\vec{s}$ is sorted,
$\partial\vec{s}$ contains only elements of form $t_{j+1}-t_j$, which is exactly
$\partial\vec{t}$.

\textbf{$\Leftarrow$ direction.} Given first condition, we know that $\vec{z}
\subseteq \vec{s}$. Thus, if we show $\vec{s} \subseteq \vec{t}$, we are done.
Suppose, on the contrary, that $\vec{s} \not\subseteq \vec{t}$. Without loss of
generality, assume thus that there is some $h \in \vec{s}$ such that
$t_{i}<h<t_{i+1}$ for some $i \in [d]$. Then, $h-t_i,t_{i+1}-h \in
\partial\vec{s}$ but clearly $h-t_i,t_{i+1}-h \not\in \partial\vec{t}$ since $h
\not\in \vec{t}$. Contradiction. Thus $\vec{z} \subseteq \vec{s} \subseteq
\vec{t}$.

\begin{example}
    Again, consider the following lookup table and (sorted) witness elements:
    \begin{equation*}
        \vec{t} = \{1,4,8\}, \quad \vec{z} = \{1,1,4,8,8,8\}, \quad \vec{z}^{\,*} = \{1,5,5,5,8,8\}.
    \end{equation*}
    Now consider the sorted multisets of concatenations:
    \begin{equation*}
        \vec{s} = \{1,1,1,4,4,8,8,8,8\}, \quad \vec{s}^{\,*} = \{1,1,4,5,5,5,8,8,8\}
    \end{equation*}

    Let us check whether both conditions are satisfied. The first condition is satisfied for 
    both $\vec{s}$ and $\vec{s}^{\,*}$ due to construction. The second condition, though, is 
    satisfied only for $\vec{s}$. Indeed:
    \begin{equation*}
        \partial\vec{t} = \{3,4\} = \partial\vec{s} \neq \partial\vec{s}^{\,*} = \{1,3\}
    \end{equation*}
\end{example}

\textbf{Reducing two checks to one.} At this point, one can apply the
$\mathcal{P}$lon$\mathcal{K}$-based protocol to check whether $\vec{s} =
(\vec{z}, \vec{t})$ and $\partial\vec{s} = \partial\vec{t}$. However,
\textit{plookup} allows to use just a single one. Here how it works.

\begin{definition}
    The \textbf{randomized difference set} with randomness $\beta \in
    \mathbb{F}^{\times}$ of multiset $\vec{s}=\{s_i\}_{i \in [n]}$, further
    denoted as $\partial_{\beta}\vec{s}$, is defined as follows:
    \begin{equation*}
        \partial_{\beta}\vec{s} \triangleq \{s_i+\beta s_{i+1}\}_{i \in [n-1]}.
    \end{equation*}
\end{definition}

\begin{remark}
    Note that $\partial_{-1}\vec{s}$ corresponds to the difference set
    of $\vec{s}$ (albeit with the negative sign).
\end{remark}

Now, the verifier chooses a random $\beta \leftarrowS \mathbb{F}^{\times}$ and
it suffices to impose a single multiset check between $\partial_{\beta}\vec{s}$
and $((1+\beta)\vec{z},\partial_{\beta}\vec{t})$. We provide the following
proposition (without proof).
\begin{proposition}
    $\vec{z} \subseteq \vec{t}$ if and only if $\partial_{\beta}\vec{s} =
    ((1+\beta)\vec{z},\partial_{\beta}\vec{t})$ for $\beta \leftarrowS
    \mathbb{F}^{\times}$.
\end{proposition}

\textbf{Intuition.} Such proposition came out of nowhere, so let us explain why
(intuitively) it works. First, let us see how $\partial_{\beta}\vec{s}$ for
$\vec{z} \subseteq \vec{t}$ looks like. Compared to $\partial\vec{s}$, two same
consecutive elements $s_i,s_{i+1}$ with $s_i=s_{i+1}$ don't get zeroed: instead,
one gets $(1+\beta)s_i$. So at least we know where $1+\beta$ coefficient comes
from! Other than same elements, in $\partial_{\beta}\vec{s}$, one gets
coefficients of form $t_i+\beta t_{i+1}$, which is exactly
$\partial_{\beta}\vec{t}$ by definition. This way $\partial_{\beta}\vec{s}$ consists of:
\begin{itemize}
    \item Expressions of form $(1+\beta)z_i$, which together form $(1+\beta)\vec{z}$.
    \item Expressions of form $t_i+\beta t_{i+1}$, which together form
    $\partial_{\beta}\vec{t}$.
\end{itemize}
Thus we naturally require $\partial_{\beta}\vec{s} =
((1+\beta)\vec{z},\partial_{\beta}\vec{t})$. Why this condition is sufficient is
less obvious, so we leave it without proof.

\begin{example}
    Once again, consider the following lookup table and (sorted) witness elements:
    \begin{equation*}
        \vec{t} = \{1,4,8\}, \quad \vec{z} = \{1,1,4,8,8,8\}.
    \end{equation*}
    As previously shown, the sorted concatenation is $\vec{s}=\{1,1,1,4,4,8,8,8,8\}$. Now, suppose 
    we have sampled some random $\beta \leftarrowS \mathbb{F}^{\times}$. Then we have:
    \begin{equation*}
        \partial_{\beta}\vec{s} = \{1+\beta,1+\beta,1+4\beta,(1+\beta)4,4+8\beta,(1+\beta)8,(1+\beta)8,(1+\beta)8\}
    \end{equation*}
    Now note that for $(1+\beta)\vec{z}$ and $\partial_{\beta}\vec{t}$ we have:
    \begin{equation*}
        (1+\beta)\vec{z} = \{1+\beta,1+\beta,(1+\beta)4,(1+\beta)8,(1+\beta)8,(1+\beta)8\}, \quad \partial_{\beta}\vec{t} = \{1+4\beta,4+8\beta\}
    \end{equation*}
    Clearly, $\partial_{\beta}\vec{s} = ((1+\beta)\vec{z},\partial_{\beta}\vec{t})$.
\end{example}

\subsubsection{plookup Precise Scheme}

Fix integers $n,d$ --- witness size and lookup table size. Given witness
$\boldsymbol{z} \in \mathbb{F}^n$ and lookup table $\boldsymbol{t} \in
\mathbb{F}^d$, we want to ensure $\boldsymbol{z} \subseteq \boldsymbol{t}$. Let
$\Omega = \{\omega^j\}_{j \in [n]} \leq \mathbb{F}^{\times}$ be a multiplicative
field subgroup. We say $\boldsymbol{z} \subseteq \boldsymbol{t}$ is
\textit{sorted by $\boldsymbol{t}$} when values appear in the same order in
$\boldsymbol{z}$ as they do in $\boldsymbol{t}$. 

Now, given $\boldsymbol{t} \in \mathbb{F}^d$, $\boldsymbol{z} \in \mathbb{F}^n$,
and $\boldsymbol{s} \in \mathbb{F}^{n+d}$, define bi-variate polynomials $Z$ and
$T$ as follows:
\begin{gather*}
    Z(\beta,\gamma) \triangleq (1+\beta)^n \prod_{i \in [n]}(\gamma+z_i)\prod_{i \in [d-1]}(\gamma(1+\beta)+t_i+\beta t_{i+1}) \\
    T(\beta,\gamma) \triangleq \prod_{i \in [n+d-1]} (\gamma(1+\beta)+s_i+\beta s_{i+1})
\end{gather*}

Why on Earth do we need these two polynomials? Consider the following
proposition:

\begin{proposition}
    $Z \equiv T$ if and only if $\boldsymbol{z} \subseteq \boldsymbol{t}$ and
    $\boldsymbol{s}$ is $(\boldsymbol{z},\boldsymbol{t})$ sorted by
    $\boldsymbol{t}$.
\end{proposition}

Now, this fact is completely unobvious, and you can see the proof in the
\href{https://eprint.iacr.org/2020/315.pdf}{original plookup paper}. This
motivates us to formulate the following protocol. Without loss of generality,
assume $d=n+1$ (if $d \leq n$, pad $\boldsymbol{t}$ with $n-d+1$ repetitions of
the last element).

\textit{Preprocessing.} Compute the polynomial $t(X)$ that encodes the values of
lookup table $\boldsymbol{t} \in \mathbb{F}^{n+1}$ over the domain $\Omega$:
that is, $t(\omega^j)=t_j$. 

\textit{Inputs.} Polynomial $z(X)$, encoding vector $\boldsymbol{z} \in
\mathbb{F}^n$ over $\Omega$: that is, $z(\omega^j) = z_j$.

\textit{Protocol.} 
\begin{enumerate}
    \item $\mathcal{P}$ computes $\boldsymbol{s} = (\boldsymbol{z},
    \boldsymbol{t}) \in \mathbb{F}^{2n+1}$ sorted by $\boldsymbol{t}$.
    Interpolate two polynomials $h_1,h_2 \in \mathbb{F}^{<(n+1)}[X]$ as follows:
    $h_1(\omega^j)=s_j$ for $j \in [n+1]$ and $h_2(\omega^j) = s_{n+j}$ for $j
    \in [n+1]$. That is, $h_1$ interpolates first $n+1$ elements of
    $\boldsymbol{s}$, while $h_2$ the remaining $n$.
    \item $\mathcal{P}$ sends commitments of $h_1$ and $h_2$ to verifier
    $\mathcal{V}$ (namely, so that $\mathcal{V}$ will have an oracle access to
    both $h_1$ and $h_2$).
    \item $\mathcal{V}$ picks random $\beta,\gamma \leftarrowS \mathbb{F}$ and
    sends them to $\mathcal{P}$.
    \item $\mathcal{P}$ computes a polynomial $F \in \mathbb{F}^{<(n+1)}[X]$
    that aggregate the value $Z(\beta,\gamma)/T(\beta,\gamma)$ where $Z,T$ are
    as described above. Specifically, let:
    \begin{enumerate}
        \item $F(\omega) = 1$.
        \item For each $2 \leq j \leq n$, compute:
        \begin{equation*}
            F(\omega^i) = \frac{(1+\beta)^{i-1}\prod_{j<i}(\gamma+z_j)\prod_{1\leq j < i}(\gamma(1+\beta)+t_j+\beta t_{j+1})}{\prod_{1 \leq j < i}(\gamma(1+\beta)+s_j+\beta s_{j+1})(\gamma (1+\beta)+s_{n+j}+\beta s_{n+j+1})}
        \end{equation*}
        \item $F(\omega^{n+1})=1$.
    \end{enumerate}
    \item $\mathcal{P}$ sends commitment of $F$ to $\mathcal{V}$.
    \item $\mathcal{V}$ checks that $F$ is of the form described above, and that
    $F(\omega^{n+1})=1$. More precisely, $\mathcal{V}$ checks the following
    identities for each $\omega^j \in \Omega$:
    \begin{enumerate}
        \item $L_1(\omega^j)(F(\omega^j)-1)=0$.
        \item
        $(\omega^j-\omega^{n+1})F(\omega^j)(1+\beta)(\gamma+z(\omega^j))(\gamma(1+\beta)+t(\omega^j)+\beta
        t(\omega \cdot \omega^j)) = (\omega^j-\omega^{n+1})F(\omega \cdot
        \omega^j)(\gamma(1+\beta)+h_1(\omega^j)+\beta h_1(\omega \cdot
        \omega^j))(\gamma(1+\beta)+h_2(\omega^j)+\beta h_2(\omega \cdot
        \omega^j))$.
        \item $L_{n+1}(\omega^j)(h_1(\omega^j)-h_2(\omega \cdot \omega^{j}))=0$.
        \item $L_{n+1}(\omega^j)(F(\omega^j)-1)=0$.
    \end{enumerate}
    Recall that $L_j(X)$ is the Lagrange basis polynomial: that is, $L_j(\omega^i)=\delta_{i,j}$.
\end{enumerate} 

\textbf{Intuition.} Four checks that $\mathcal{V}$ imposes are not very 
obvious (especially the second one). So here is the brief explanation of their meaning:
\begin{enumerate}
    \item First equation checks whether $F(\omega)=1$: note that since
    $L_1(\omega^j)$ is non-zero only for $j=1$, the equation essentially reduces
    to $F(\omega)-1=0$.
    \item This is the main verification that is needed to check whether
    $Z(\beta,\gamma) \equiv T(\beta,\gamma)$. However, since both $Z$ and $T$
    are high-degree polynomials, we cannot directly check their equivalence. For 
    that reason, we introduce the ``accumulator'' polynomial $F$ that should satisfy 
    $F(\omega) = F(\omega^{n+1}) = 1$ and recursive condition
    \begin{equation*}
        \frac{F(\omega^{j+1})}{F(\omega^j)} = \frac{(1+\beta)(\gamma+z_j)(\gamma(1+\beta)+t_j+\beta t_{j+1})}{(\gamma(1+\beta)+s_j+\beta s_{j+1})(\gamma(1+\beta) + s_{n+j} + \beta s_{n+j+1})}
    \end{equation*}
    Note that given $F(\omega)=1$ and this recursive relation, by taking the
    product of both sides from $1$ to $n$ one obtains
    $F(\omega^{n+1})=Z(\beta,\gamma)/T(\beta,\gamma)$ which is one, according to the
    previously specified lemma.
    \item Since $L_{n+1}(\omega^j)$ is non-zero only for $j=n+1$, this check
    verifies that $h_1(\omega^{n+1}) = h_2(\omega^{n+2})$. This is the
    consistency check for $h_1$ and $h_2$ that checks whether $h_1$ and $h_2$
    are properly ``glued'' together to encode $\boldsymbol{s}$.
    \item Finally, the last check verifies whether $F(\omega^{n+1})=1$.
\end{enumerate}

\begin{lemma}
    Soundness of the above protocol is at least $1-(5n+4)/|\mathbb{F}|$.
\end{lemma}

\subsection{Logup}

While \textit{plookup} is the Poly-IOP-based protocol, \textit{logup} allows to
prove the lookup inclusion using SumCheck and rather more algebraic (compared to
plookup) approach. Let's see how it works.

\subsubsection{Preliminaries}

Compared to previously discussed Sum-Check protocol, to stick with the original
logup notation, we consider a slighly different hypercube: instead of boolean
version $\{0,1\}^n$, we consider the hypercube $\{-1,1\}^n=\{\pm 1\}^n$, which
we denote by $\mathcal{Q}_n$ for short. Note that the SumCheck protocol is
exactly the same as for the boolean case, but we need to use ``$-1$'' instead of
``$0$''. Besides, the multilinear lagrange basis polynomials have a much 
nicer form:
\begin{equation*}
    \mathsf{eq}(\mathbf{x};\mathbf{y}) = \frac{1}{2^n}\prod_{j=1}^n (1+x_jy_j)
\end{equation*}

\subsubsection{Logarithmic Derivative}

One of the ideas of the logup protocol is to reduce the polynomial check 
to the fractional one. To show this transition, we need a couple of formalities 
before we proceed.

\begin{definition}
    Given a polynomial $q(X) := \sum_{j=0}^d q_jX^j \in \mathbb{F}[X]$, its
    \textbf{formal derivative}, in a similar fashion to calculus, is given by
    $q'(X) \triangleq \sum_{j=1}^{d}jq_jX^{j-1}$.
\end{definition}

\begin{example}
    For instance, given $q(X)=1+2X+3X^2$, the formal derivative is $q'(X)=2+6X$.
\end{example}

It can be shown that the formal derivative obeys all the usual rules taught 
in calculus. Now, one more definition.
\begin{definition}
    For a function $q(X)/r(X)$ from the rational function field $\mathbb{F}(X)$, the 
    formal derivative is defined as:
    \begin{equation*}
        \left(\frac{q(X)}{r(X)}\right)' \triangleq \frac{q'(X)r(X)-q(X)r'(X)}{r(X)^2}
    \end{equation*}
\end{definition}

For the further derivation, we will be interested in particular cases when
derivatives turn out to be zero: that is $q'(X)=0$ for the polynomial case or
$\left(\frac{q(X)}{r(X)}\right)'=0$ for a rational one. If polynomials are given
in the real field, then these conditions immediately imply $q(X) \equiv
\text{const}$ for polynomial case and $q(X)=cr(X)$ with $c=\text{const}$ for the
rational case. Surprisingly, when considering the field $\mathbb{F}$ of finite
characteristic $p$, these results are not generally true. Fortunately for us,
they turn out to be false only for extreme case when degrees of polynomials are
more than $p$, which, as you can imagine, never happens in practice. So we give
the following lemma.
\begin{lemma}
    Assume $q(X)$ and $r(X)$ are polynomials of degree less than $p$. Then:
    \begin{itemize}
        \item If $q'(X)=0$, then $q$ is constant.
        \item If $(q(X)/r(X))'=0$, then $q(X)/r(X)=c$ for some $c \in
        \mathbb{F}$.
    \end{itemize}
\end{lemma}

Finally, the logarithmic derivative.
\begin{definition}
    The \textbf{logarithmic derivative} of $q(X) \in \mathbb{F}[X]$ is given by
    the rational function $q'(X)/q(X)$.
\end{definition}

\begin{remark}
    The name ``logarithmic'' comes from the fact that in calculus, for a
    function $f: \mathbb{R} \to \mathbb{R}$, the derivative of $\log f(x)$ is
    given by $f'(x)/f(x)$ by the chain rule.
\end{remark}

One reason why logarithmic derivatives are attractive for our application is the
following: the logarithmic derivative of $q(X)r(X)$ is the sum of logarithmic
derivatives of $q(X)$ and $r(X)$. Indeed, note that the logarithmic derivative 
of $q(X)r(X)$ is nothing but:
\begin{equation*}
    \frac{(q(X)r(X))'}{q(X)r(X)} = \frac{q'(X)r(X) + q(X)r'(X)}{q(X)r(X)} = \frac{q'(X)}{q(X)} + \frac{r'(X)}{r(X)}
\end{equation*}

In particular, logarithmic derivative of $\prod_{j=1}^n (X+z_j)$ is given by
$\sum_{j=1}^n \frac{1}{X+z_j}$. From all previously mentioned facts we can 
finally formulate some useful facts about lookup checks. Let us start 
from the simple observation.

\begin{lemma}
    Let $\vec{a} = \{a_i\}_{i \in [n]}$ and $\vec{b} = \{b_i\}_{i \in [n]}$ be
    two sequences of elements from $\mathbb{F}$. To verify with overwhelming
    probability whether two multisets are equal, it suffices to check
    \begin{equation*}
        \sum_{j=1}^n \frac{1}{\gamma + a_j} = \sum_{j=1}^n \frac{1}{\gamma+b_j}
    \end{equation*}
    for randomly chosen $\gamma \leftarrowS \mathbb{F}$.
\end{lemma}

\textbf{Proof.} Recall that two multisets are equal with overwhelming
probability if for randomly selected $\gamma \leftarrowS \mathbb{F}$, we have
$\prod_{j=1}^n (a_j+\gamma) = \prod_{j=1}^n (b_j+\gamma)$. Thus it suffices to
show that imposing this product equality is equivalent to the statement in the
lemma. Notice that if products are indeed the same, so are their logarithmic
derivatives (with respect to $\gamma$), so we immediately get $\sum_{j=1}^n
\frac{1}{\gamma + a_j} = \sum_{j=1}^n \frac{1}{\gamma+b_j}$. To show the
opposite direction, assume that logarithmic derivatives of $q_{\vec{a}}(X) =
\prod_{j=1}^n (X+a_j)$ and $q_{\vec{b}}(X) = \prod_{j=1}^n (X+b_j)$ are the same,
so $\frac{q_{\vec{a}}'(X)}{q_{\vec{a}}(X)} = \frac{q_{\vec{b}}'(X)}{q_{\vec{b}}(X)}$. Then:
\begin{equation*}
    \left(\frac{q_{\vec{a}}(X)}{q_{\vec{b}}(X)}\right)' = \frac{q_{\vec{a}}'(X)q_{\vec{b}}(X)-q_{\vec{a}}(X)q_{\vec{b}}'(X)}{q_{\vec{b}}(X)^2}=0
\end{equation*}
Thus $\frac{q_{\vec{a}}(X)}{q_{\vec{b}}(X)}=c$ for constant $c \in \mathbb{F}$.
Since the leading coefficients of both $q_{\vec{a}}(X)$ and $q_{\vec{b}}(X)$ are
$1$, we conclude that $c=1$, which immediately implies $\prod_{j=1}^n
(a_j+\gamma) = \prod_{j=1}^n (b_j+\gamma)$.

The consequence of this fact is the following.

\begin{lemma}
    Given two sequences of elements $\{t_i\}_{i \in [d]}$ and $\{z_i\}_{i \in
    [n]}$, the inclusion check $\{z_i\}_{i \in [n]} \subseteq \{t_i\}_{i \in
    [d]}$ is satisfied if and only if there exist the set of multiplicities
    $\{\mu_i\}_{i \in [d]}$ where $\mu_i = \#\{j \in [n]: z_j = t_i\}$ such that:
    \begin{equation*}
        \sum_{i \in [n]} \frac{1}{X+z_i} = \sum_{i \in [d]} \frac{\mu_i}{X+t_i}
    \end{equation*}

    In particular, checking such equality at random point from $\mathbb{F}$
    results in the soundness error of up to $(n+d)/|\mathbb{F}|$, which becomes
    negligible for fairly large $|\mathbb{F}|$.
\end{lemma}

\subsubsection{Applying SumCheck}

Now we need to build the SumCheck equation. Instead of considering the inclusion
check for scalars, we consider the inclusion check for \textit{multivariate
polynomials}. Notice that the lookup table $\vec{t} = \{t_j\}_{j \in [d]}$ for
$d=2^v$ can be viewed as a function $t: \mathcal{Q}_v \to \mathbb{F}$. However,
here is the issue: our witness vector $\vec{z} = \{z_i\}_{i \in [n]}$ typically
has a significantly larger size than $\vec{t}$, so we cannot view $\vec{z}$ as a
function $\mathcal{Q}_v \to \mathbb{F}$ since the input domain size is too small
(also quite obvious remark --- we cannot have two different hypercubes for the
sumcheck protocol). For that reason, we consider a list of $m := \lceil n/d
\rceil$ functions $z_1,\dots,z_m: \mathcal{Q}_v \to \mathbb{F}$ that encode all
$n$ values of $\{z_i\}_{i \in [n]}$.

This way, we want to check whether $\bigcup_{i \in
[m]}\{z_i(\mathbf{x})\}_{\mathbf{x} \in \mathcal{Q}_v} \subseteq
\{t(\mathbf{x})\}_{\mathbf{x} \in \mathcal{Q}_v}$. This way, we can rewrite our
condition $\sum_{i \in [n]} \frac{1}{X+z_i} = \sum_{i \in [d]}
\frac{\mu_i}{X+t_i}$ in the ``sumcheck'' world as follows:
\begin{equation*}
    \sum_{\mathbf{x} \in \mathcal{Q}_v}\sum_{i \in [m]} \frac{1}{\gamma + z_i(\mathbf{x})} = \sum_{\mathbf{x} \in \mathcal{Q}_v} \frac{\mu(\mathbf{x})}{\gamma + t(\mathbf{x})}
\end{equation*}

Here, $\mu(\mathbf{x})$ (if it is injective, which is typically the case) gives the
number of occurancies of $t(\mathbf{x})$ in $z_1,\dots,z_m$ altogether, i.e.
$\mu(\mathbf{x}) = \sum_{i \in [m]}\#\{\mathbf{y} \in \mathcal{Q}_v:
z_i(\mathbf{y}) = t(\mathbf{x})\}$.

Now, the idea of logup is to sample random $\gamma \leftarrowS \mathbb{F}$ and 
apply the SumCheck on the function:
\begin{equation*}
    \zeta(\mathbf{x}) = \sum_{i \in [m]} \frac{1}{\gamma + \widetilde{z}_i(\mathbf{x})} - \frac{\widetilde{\mu}(\mathbf{x})}{\gamma + \widetilde{t}(\mathbf{x})},
\end{equation*}

where by $\widetilde{t}_i$ we denoted the multilinear extension, that is:
\begin{equation*}
\widetilde{t}(\mathbf{x}) = \sum_{\mathbf{y} \in \mathcal{Q}_v}
t(\mathbf{y})\mathsf{eq}(\mathbf{x}; \mathbf{y}) = \frac{1}{2^v}\sum_{\mathbf{y} \in \mathcal{Q}_v}t(\mathbf{y})\prod_{j=1}^v(1+x_jy_j) 
\end{equation*}

The only issue left is that sumcheck protocol cannot work with rational
functions, which is the case here for $\zeta(\mathbf{x})$. For that reason,
roughly, the prover $\mathcal{P}$ will divide the sum of $m$ terms into $\ell$
chunks and provide multilinear helper functions for each such sum. Note that 
$\ell$ is chosen depending on the chosen polynomial commitment scheme. 

Namely, suppose we split $[m] = \bigcup_{j \in [k]}\mathcal{I}_j$ into $k=\lceil
m/\ell \rceil$ subintervals. Suppose 
\begin{equation*}
    \zeta_j(\mathbf{x}) = \sum_{i \in \mathcal{I}_j} \frac{\mu_i(\mathbf{x})}{\phi_i(\mathbf{x})}, \quad j \in [k],
\end{equation*}
is the respective partial sum of consecutive terms in the overall expression
$\frac{\mu(\mathbf{x})}{\gamma+\widetilde{t}(\mathbf{x})} - \sum_{i \in
[m]}\frac{1}{\gamma+\widetilde{z}_i(\mathbf{x})}$. That is, we used the notation
$\mu_0(\mathbf{x}) = \mu(\mathbf{x})$ and $\mu_i(\mathbf{x}) \equiv -1$ for
$i>0$. Similarly, $\phi_0(\mathbf{x}) = \gamma+\widetilde{t}(\mathbf{x})$ and
$\phi_i(\mathbf{x}) = \gamma +\widetilde{z}_{i-1}(\mathbf{x})$ for $i>0$.

Then, the prover provides oracles for $\{\zeta_i\}_{i \in [k]}$ subject to
$\sum_{\mathbf{x} \in \mathcal{Q}_v}\sum_{i \in [k]}\zeta_i(\mathbf{x}) = 0$
and the domain identities:
\begin{equation*}
    \zeta_j(\mathbf{x}) \prod_{i \in \mathcal{I}_j} \phi_i(\mathbf{x}) = \sum_{i \in \mathcal{I}_j}\mu_i(\mathbf{x})\prod_{k \in \mathcal{I}_j\setminus \{i\}} \phi_j(\mathbf{x})
\end{equation*}

These identities are finally checked using SumCheck protocol and then combined
into a single one using random scalars $\lambda_0,\dots,\lambda_{k-1}
\leftarrowS \mathbb{F}$.

\subsubsection{Precise Scheme}

Suppose the prover $\mathcal{P}$ wants to convince the verifier $\mathcal{V}$
that $\{z_i\}_{i \in [n]} \subseteq \{t_j\}_{j \in [d]}$. Assume $d=2^v$. 

\textit{Preprocessing stage.} Compute multilinear extension $\widetilde{t}:
\mathbb{F}^v \to \mathbb{F}$ that encodes the lookup table $\{t_j\}_{j \in [d]}$
and, if needed, multilinear Lagrange basis polynomials
$\{\mathsf{eq}(\mathbf{x};\mathbf{y})\}_{\mathbf{y} \in \mathcal{Q}_v}$.

\textit{Protocol.} Interaction between $\mathcal{P}$ and $\mathcal{V}$ proceeds as follows:
\begin{enumerate}
    \item $\mathcal{P}$ divides witness into $k = \lceil n/d \rceil$ pieces and
    perceives values $\{z_i\}_{i \in [n]}$ as a set of $m$ functions $z_j:
    \{\pm 1\}^v \to \mathbb{F}$. Additionally, $\mathcal{P}$ computes the
    multilinear extensions $\{\widetilde{z}_j(\mathbf{x})\}_{j \in [m]}$.
    \item $\mathcal{P}$ computes the multiplicity function multilinear extension
    $\widetilde{\mu}$ and sends the oracle access to it
    $\mathcal{O}^{\mu}(\cdot)$ to the verifier $\mathcal{V}$.
    \item $\mathcal{V}$ samples the challenge $\gamma \leftarrowS \mathbb{F}$
    and sends to $\mathcal{P}$.
    \item $\mathcal{P}$ computes $k=\lceil m/\ell \rceil$ functions
    $\{\zeta_j(\mathbf{x})\}_{j \in [k]}$ according to constraints above and
    sends oracle access
    $\mathcal{O}^{\zeta_0}(\cdot),\mathcal{O}^{\zeta_1}(\cdot),\dots,\mathcal{O}^{\zeta_{k-1}}(\cdot)$
    to the verifier $\mathcal{V}$.
    \item $\mathcal{V}$ responds with a random vector $\boldsymbol{\alpha} \leftarrowS \mathbb{F}^n$ and 
    random scalars $\lambda_0,\dots,\lambda_{k-1} \leftarrowS \mathbb{F}$. Now, both $\mathcal{P}$ and 
    $\mathcal{V}$ engage in the sumcheck protocol for:
    \begin{equation*}
        \sum_{\mathbf{x} \in \mathcal{Q}_v}G(\mathbf{x}, \mathsf{eq}(\mathbf{x};\boldsymbol{\alpha}),\mu(\mathbf{x}),\phi_0(\mathbf{x}),\dots,\phi_{m-1}(\mathbf{x}),\zeta_0(\mathbf{x}),\dots,\zeta_{k-1}(\mathbf{x})) = 0,
    \end{equation*}
    where the function $G$ is defined as follows:
    \begin{equation*}
        G(\mathbf{x},\star) = \sum_{r \in [k]} \zeta_r(\mathbf{x}) + \mathsf{eq}(\mathbf{x};\boldsymbol{\alpha})\lambda_r\left(\zeta_r\prod_{i \in \mathcal{I}_r}\phi_i(\mathbf{x}) - \sum_{i \in \mathcal{I}_r}\mu_i(\mathbf{x})\prod_{j \in \mathcal{I}_r \setminus \{i\}} \phi_j(\mathbf{x})\right)
    \end{equation*}
    During sumcheck, $\mathcal{V}$ uses oracles provided by $\mathcal{P}$.
\end{enumerate}

\begin{lemma}
    The cost for running the protocol is $\mathcal{O}(n \ell)$
    field multiplications, while the communication size is $\mathcal{O}(m/\ell)$ oracles of size
    $\mathcal{O}(d)$.
\end{lemma}
